% Lineare Funktionen – bank/lineare-funktionen/ (zusammenbau v0.4, Heft)
\einheitenkopf[t1]{Lineare Funktionen}
\setcounter{aufgabe}{0}

% Anlauf (ohne Prüfkennung)
\begin{aufgabe}{Graph zeichnen – Anlauf}
\begin{teile}
\teil Wertetabelle: Fülle die Tabelle zu $f(x) = 3x - 2$ aus.

\wertetabelle{x}{f(x)}{-1,0,1,2}
\teil Mit n und Steigungsdreieck: Zeichne die Gerade zu $f(x) = -3x + 4$.

\begin{ksys}[xmin=-4,xmax=4,ymin=-4,ymax=5] \end{ksys}
\end{teile}
\end{aufgabe}

\begin{aufgabe}{Graph zeichnen – Prüfungsaufgaben}
\begin{teile}
\teil Graph zeichnen: Zeichne den Graphen der Funktion f mit $f(x) = -4x + 3$ in das Koordinatensystem. \hfill (P10 2026 FOR)

\begin{ksys}[xmin=-4,xmax=4,ymin=-4,ymax=4] \end{ksys}
\teil Graph zeichnen: Zeichne den Graphen der Funktion f mit $f(x) = -3x - 4$ in das Koordinatensystem. \hfill (P10 2026 FOR)

\begin{ksys}[xmin=-4,xmax=4,ymin=-8,ymax=4] \end{ksys}
\teil Graph zeichnen: Zeichne den Graphen der Funktion f mit $f(x) = 5x - 3$ in das Koordinatensystem. \hfill (P10 2024 OS)

\begin{ksys}[xmin=-4,xmax=4,ymin=-4,ymax=4] \end{ksys}
\teil Graph zeichnen: Zeichne den Graphen der Funktion f mit $f(x) = 6x - 4$ in das Koordinatensystem. \hfill (P10 2024 OS)

\begin{ksys}[xmin=-4,xmax=4,ymin=-5,ymax=4] \end{ksys}
\teil Graph zeichnen: Zeichne den Graphen der Funktion f mit $f(x) = -3x + 5$ in das Koordinatensystem. \hfill (P10 2022 OS)

\begin{ksys}[xmin=-4,xmax=4,ymin=-4,ymax=6] \end{ksys}
\teil Graph zeichnen: Zeichne den Graphen der Funktion f mit $f(x) = -4x - 2$ in das Koordinatensystem. \hfill (P10 2022 OS)

\begin{ksys}[xmin=-4,xmax=4,ymin=-7,ymax=4] \end{ksys}
\teil Graph zeichnen: Zeichne den Graphen der Funktion f mit $f(x) = 2x + 5$ in das Koordinatensystem. \hfill (P10 2021 OS)

\begin{ksys}[xmin=-4,xmax=4,ymin=-4,ymax=8] \end{ksys}
\teil Graph zeichnen: Zeichne den Graphen der Funktion f mit $f(x) = 4x - 3$ in das Koordinatensystem. \hfill (P10 2021 OS)

\begin{ksys}[xmin=-4,xmax=4,ymin=-4,ymax=4] \end{ksys}
\end{teile}
\end{aufgabe}

% Anlauf (ohne Prüfkennung)
\begin{aufgabe}{Ablesen – Anlauf}
\begin{teile}
\teil n? Lies den y-Achsenabschnitt n der Geraden g ab.

\begin{ksys}[xmin=-5,xmax=5,ymin=-5,ymax=5,ablesen] \gerade{1}{3}{g} \end{ksys}
n = \leerfeld
\teil m? Lies die Steigung m der Geraden g mit einem Steigungsdreieck ab.

\begin{ksys}[xmin=-5,xmax=5,ymin=-5,ymax=5,ablesen] \gerade{-2}{1}{g} \end{ksys}
m = \leerfeld
\end{teile}
\end{aufgabe}

\begin{aufgabe}{Ablesen – Prüfungsaufgaben}
\begin{teile}
\teil Welche Gleichung gehört zu welcher Geraden? Ordne den Geraden f, g und h je eine der Gleichungen zu: $y = -3x + 3$; $y = -\frac{1}{3}x + 3$; $y = -3$; $y = 2x + 3$; $y = -3x$; $y = 2x - 3$. \hfill (P10 2019 OS) \\

\begin{ksys}[xmin=-5,xmax=5,ymin=-5,ymax=5,ablesen] \gerade{-3}{3}{f} \gerade{2}{3}{g} \gerade{0}{-3}{h} \end{ksys}
f: \leerfeld \quad g: \leerfeld \quad h: \leerfeld
\teil Welche Gleichung gehört zu welcher Geraden? Ordne den Geraden f, g und h je eine der Gleichungen zu: $y = -4x + 4$; $y = -\frac{1}{4}x + 4$; $y = -1$; $y = 2x + 4$; $y = -4x$; $y = 2x - 4$. \hfill (P10 2019 OS) \\

\begin{ksys}[xmin=-5,xmax=5,ymin=-5,ymax=5,ablesen] \gerade{-4}{4}{f} \gerade{2}{4}{g} \gerade{0}{-1}{h} \end{ksys}
f: \leerfeld \quad g: \leerfeld \quad h: \leerfeld
\teil Welche Gerade hat die Eigenschaft? Gib zu jeder Eigenschaft die passende Gerade an: Anstieg $-3$; parallel zur x-Achse. \hfill (P10 2019 OS) \\

\begin{ksys}[xmin=-5,xmax=5,ymin=-5,ymax=5,ablesen] \gerade{-3}{1}{f} \gerade{3}{1}{g} \gerade{0}{3}{h} \end{ksys}
Anstieg: \leerfeld \quad parallel zur x-Achse: \leerfeld
\teil Welche Gerade hat die Eigenschaft? Gib zu jeder Eigenschaft die passende Gerade an: Anstieg $-2$; parallel zur x-Achse. \hfill (P10 2019 OS) \\

\begin{ksys}[xmin=-5,xmax=5,ymin=-5,ymax=5,ablesen] \gerade{2}{-1}{f} \gerade{-2}{-1}{g} \gerade{0}{-3}{h} \end{ksys}
Anstieg: \leerfeld \quad parallel zur x-Achse: \leerfeld
\teil Welche Gerade fällt? Kreuze an. \hfill (P10 2016 OS) \\ \kreuz{f} \\ \kreuz{g}

\begin{ksys}[xmin=-5,xmax=5,ymin=-5,ymax=5,ablesen] \gerade{1.5}{-1}{f} \gerade{-0.5}{2}{g} \end{ksys}
\teil Welche Gerade fällt? Kreuze an. \hfill (P10 2016 OS) \\ \kreuz{f} \\ \kreuz{g}

\begin{ksys}[xmin=-5,xmax=5,ymin=-5,ymax=5,ablesen] \gerade{-1}{3}{f} \gerade{2}{-3}{g} \end{ksys}
\teil Schnittpunkt mit der y-Achse: Markiere ihn an der Geraden $y = 1{,}5x - 3$. \hfill (P10 2024 OS)

\begin{ksys}[xmin=-5,xmax=5,ymin=-5,ymax=5,ablesen] \gerade{1.5}{-3}{g} \end{ksys}
\teil Schnittpunkt mit der y-Achse: Markiere ihn an der Geraden $y = -1{,}5x + 3$. \hfill (P10 2024 OS)

\begin{ksys}[xmin=-5,xmax=5,ymin=-5,ymax=5,ablesen] \gerade{-1.5}{3}{g} \end{ksys}
\teil Welcher Graph gehört zu einer linearen Funktion? Kreuze an. \hfill (P10 2025 OS) \\ \kreuz{A} \\ \kreuz{B} \\ \kreuz{C} \\ \kreuz{D}

\begin{ksys}[xmin=-5,xmax=5,ymin=-5,ymax=5,ablesen] \gerade{-0.5}{3}{A} \funktionab{2/\x}{B}{0.4}{5} \parabel{1}{2}{-4}{C} \funktion{abs(\x+2)-4}{D} \end{ksys}
\teil Welcher Graph gehört zu einer linearen Funktion? Kreuze an. \hfill (P10 2025 OS) \\ \kreuz{A} \\ \kreuz{B} \\ \kreuz{C} \\ \kreuz{D}

\begin{ksys}[xmin=-5,xmax=5,ymin=-5,ymax=5,ablesen] \funktion{abs(\x-2)+1}{A} \parabel{-1}{-2}{4}{B} \gerade{1.5}{-3}{C} \funktionab{-3/\x}{D}{0.6}{5} \end{ksys}
\end{teile}
\end{aufgabe}

% Anlauf (ohne Prüfkennung)
\begin{aufgabe}{Funktionswert – Anlauf}
\begin{teile}
\teil $f(4)$? Berechne den Funktionswert von $f(x) = 2x + 5$ an der Stelle $x = 4$. f(4) = \leerfeld
\teil $f(-3)$? Berechne den Funktionswert von $f(x) = 4x - 1$ an der Stelle $x = -3$. f(-3) = \leerfeld
\end{teile}
\end{aufgabe}

\begin{aufgabe}{Funktionswert – Prüfungsaufgaben}
\begin{teile}
\teil Graph und Nullstelle: Zeichne den Graphen der Funktion f mit $f(x) = -4x + 6$ und berechne die Nullstelle von f. \hfill (P10 2022 OS) \\

\begin{ksys}[xmin=-4,xmax=4,ymin=-4,ymax=7] \end{ksys}
x = \leerfeld
\teil Graph und Nullstelle: Zeichne den Graphen der Funktion f mit $f(x) = -6x + 3$ und berechne die Nullstelle von f. \hfill (P10 2022 OS) \\

\begin{ksys}[xmin=-4,xmax=4,ymin=-4,ymax=4] \end{ksys}
x = \leerfeld
\teil Wann ist der Tank leer? Ein Wassertank enthält 60 Liter. Die Gleichung $y = -0{,}4x + 60$ gibt an, wie viele Liter y nach x Minuten noch im Tank sind. Berechne, nach wie vielen Minuten der Tank leer ist. \hfill (P10 2021 OS) \\ \leerfeld[min]
\teil Wann ist der Akku leer? Ein voll geladener Handyakku verliert beim Spielen 0{,}8 Prozentpunkte je Minute. Die Gleichung $y = -0{,}8x + 100$ gibt den Ladestand y in Prozent nach x Minuten an. Berechne, nach wie vielen Minuten der Akku leer ist. \hfill (P10 2021 OS) \\ \leerfeld[min]
\teil Nullstelle ablesen: Gib die Nullstelle der Funktion g an. \hfill (P10 2019 OS) \\

\begin{ksys}[xmin=-5,xmax=5,ymin=-5,ymax=5,ablesen] \gerade{1.5}{3}{g} \end{ksys}
x = \leerfeld
\teil Nullstelle ablesen: Gib die Nullstelle der Funktion g an. \hfill (P10 2019 OS) \\

\begin{ksys}[xmin=-5,xmax=5,ymin=-5,ymax=5,ablesen] \gerade{-0.5}{2}{g} \end{ksys}
x = \leerfeld
\teil Gerade und Parabel: Zeichne die Gerade f mit $f(x) = 3x - 3$ in das Koordinatensystem, gib die Koordinaten des sichtbaren Schnittpunkts von f mit der Parabel p an und entscheide, ob sich f und p in zwei Punkten schneiden. \hfill (P10 2023 OS) \\

\begin{ksys}[xmin=-5,xmax=5,ymin=-6,ymax=8] \parabel{-1}{1}{4}{p} \end{ksys}
S( \leerfeld | \leerfeld ) \quad \janein
\teil Gerade und Parabel: Zeichne die Gerade f mit $f(x) = 4x + 2$ in das Koordinatensystem, gib die Koordinaten des sichtbaren Schnittpunkts von f mit der Parabel p an und entscheide, ob sich f und p in zwei Punkten schneiden. \hfill (P10 2023 OS) \\

\begin{ksys}[xmin=-5,xmax=5,ymin=-6,ymax=8] \parabel{-1}{0}{7}{p} \end{ksys}
S( \leerfeld | \leerfeld ) \quad \janein
\teil y-Koordinate: Der Punkt A liegt auf der Geraden $y = \frac{1}{3}x - 2$ und hat die x-Koordinate $-9$. Berechne die y-Koordinate von A. \hfill (P10 2017 OS) \\ y = \leerfeld
\teil y-Koordinate: Der Punkt A liegt auf der Geraden $y = \frac{3}{4}x + 2$ und hat die x-Koordinate $-8$. Berechne die y-Koordinate von A. \hfill (P10 2017 OS) \\ y = \leerfeld
\teil Liegt P auf der Geraden? Prüfe rechnerisch, ob $P(-4|5)$ auf der Geraden $f(x) = -\frac{1}{2}x + 3$ liegt. \hfill (P10 2026 FOR) \\ \janein
\teil Liegt P auf der Geraden? Prüfe rechnerisch, ob $P(-2|4)$ auf der Geraden $f(x) = -\frac{3}{2}x - 1$ liegt. \hfill (P10 2026 FOR) \\ \janein
\teil Welche zwei Aussagen stimmen? Gegeben ist $f(x) = 2x + 6$. Kreuze die beiden richtigen Aussagen an. \hfill (P10 2021 OS) \\ \kreuz{$Q(6|0)$ liegt auf der Geraden f.} \\ \kreuz{Die Gerade f fällt.} \\ \kreuz{Die Gerade f schneidet die y-Achse in $P(0|6)$.} \\ \kreuz{Die Gerade f ist parallel zur Geraden $y = 2x$.}
\teil Welche zwei Aussagen stimmen? Gegeben ist $f(x) = -3x + 9$. Kreuze die beiden richtigen Aussagen an. \hfill (P10 2021 OS) \\ \kreuz{$Q(-3|0)$ liegt auf der Geraden f.} \\ \kreuz{Die Gerade f fällt.} \\ \kreuz{Die Gerade f schneidet die y-Achse in $P(0|-3)$.} \\ \kreuz{Der Punkt $R(2|3)$ liegt auf der Geraden f.}
\teil Welcher Punkt liegt nicht auf der Geraden? Gegeben ist $f(x) = -6x + 2$. Kreuze den Punkt an, der nicht auf der Geraden f liegt. \hfill (P10 2023 OS) \\ \kreuz{$(-1|8)$} \\ \kreuz{$(2|-10)$} \\ \kreuz{$(-3|16)$}
\teil Welcher Punkt liegt nicht auf der Geraden? Gegeben ist $f(x) = -5x + 4$. Kreuze den Punkt an, der nicht auf der Geraden f liegt. \hfill (P10 2023 OS) \\ \kreuz{$(1|-1)$} \\ \kreuz{$(-2|-6)$} \\ \kreuz{$(3|-11)$}
\end{teile}
\end{aufgabe}

% Anlauf (ohne Prüfkennung)
\begin{aufgabe}{Gleichung bestimmen – Anlauf}
\begin{teile}
\teil Gleichung? Bestimme die Gleichung der Geraden g aus dem Graphen.

\begin{ksys}[xmin=-5,xmax=5,ymin=-5,ymax=5,ablesen] \gerade{2}{-3}{g} \end{ksys}
f(x) = \leerfeld
\teil Gleichung aus zwei Punkten: Bestimme die Gleichung der Geraden durch $A(1|6)$ und $B(3|0)$. f(x) = \leerfeld
\end{teile}
\end{aufgabe}

\begin{aufgabe}{Gleichung bestimmen – Prüfungsaufgaben}
\begin{teile}
\teil Gerade durch A und B: Zeichne die Gerade f durch $A(-2|8)$ und $B(4|-1)$, entscheide für jede Aussage, ob sie wahr oder falsch ist, und gib eine Gleichung von f an. \hfill (P10 2025 OS) \\ \kreuz{wahr}\kreuz{falsch} f verläuft monoton steigend. \\ \kreuz{wahr}\kreuz{falsch} f schneidet die y-Achse in $(0|5)$.

\begin{ksys}[xmin=-4,xmax=5,ymin=-4,ymax=9] \end{ksys}
f(x) = \leerfeld
\teil Gerade durch A und B: Zeichne die Gerade f durch $A(-4|-5)$ und $B(2|-2)$, entscheide für jede Aussage, ob sie wahr oder falsch ist, und gib eine Gleichung von f an. \hfill (P10 2025 OS) \\ \kreuz{wahr}\kreuz{falsch} f verläuft monoton fallend. \\ \kreuz{wahr}\kreuz{falsch} f schneidet die y-Achse in $(0|-3)$.

\begin{ksys}[xmin=-5,xmax=4,ymin=-6,ymax=4] \end{ksys}
f(x) = \leerfeld
\steil Gleichung im Sachzusammenhang: Ein Heißluftballon sinkt gleichmäßig. 10 Minuten nach Beginn der Messung ist er in 800 m Höhe, nach 40 Minuten in 440 m Höhe. Bestimme die Gleichung $h(t)$ für die Höhe h in m nach t Minuten. \hfill (P10 2015 OS) \\ h(t) = \leerfeld
\steil Gleichung im Sachzusammenhang: Ein Regenfass wird gleichmäßig gefüllt. 4 Minuten nach Beginn der Messung sind 260 Liter darin, nach 10 Minuten 500 Liter. Bestimme die Gleichung $V(t)$ für die Wassermenge V in Litern nach t Minuten. \hfill (P10 2015 OS) \\ V(t) = \leerfeld
\teil Nachweis: Zeichne die Gerade g durch $C(-3|-4)$ und $D(3|0)$ in ein Koordinatensystem. Weise nach, dass g die Gleichung $y = \frac{2}{3}x - 2$ hat. \hfill (P10 2017 OS)

\begin{ksys}[xmin=-4,xmax=4,ymin=-5,ymax=4] \end{ksys}
\teil Nachweis: Zeichne die Gerade g durch $E(-2|5)$ und $F(4|2)$ in ein Koordinatensystem. Weise nach, dass g die Gleichung $y = -\frac{1}{2}x + 4$ hat. \hfill (P10 2017 OS)

\begin{ksys}[xmin=-4,xmax=5,ymin=-4,ymax=6] \end{ksys}
\end{teile}
\end{aufgabe}

% Anlauf (ohne Prüfkennung)
\begin{aufgabe}{Anwendung – Anlauf}
\begin{teile}
\teil Markiere: Unterstreiche den Anfangswert, kreise die Änderung je Einheit ein. Ein Handwerker berechnet 45 € für die Anfahrt und 38 € je Arbeitsstunde.
\teil Welche Gleichung passt? Ein Fahrradverleih nimmt 8 € Grundgebühr und 3 € je Stunde. y ist der Preis in €, x die Anzahl der Einheiten. Kreuze an. \\ \kreuz{$y = 8x + 3$} \\ \kreuz{$y = 3x + 8$} \\ \kreuz{$y = 11x$} \\ \kreuz{$y = 3x - 8$}
\teil Stelle eine Gleichung auf: Ein Kino-Abo kostet einmalig 10 € und 6 € je Film. y sind die Kosten in €, x die Anzahl der Filme. y = \leerfeld
\end{teile}
\end{aufgabe}

\begin{aufgabe}{Anwendung – Prüfungsaufgaben}
\begin{teile}
\steil Welches Mietauto? Angebot 1: 32 € Miete pro Tag, 300 km insgesamt frei, jeder weitere km 0{,}30 €. Angebot 2: einmalig 95 € für eine Woche, jeder km 0{,}12 €. Jana mietet 4 Tage und fährt 420 km. Berechne die Gesamtkosten beider Angebote, entscheide, welches günstiger ist, und stelle für Angebot 2 eine Gleichung für die Kosten y (in €) bei x km für eine Woche auf. \hfill (P10 2023 OS) \\ Angebot 1: \leerfeld[€] \quad Angebot 2: \leerfeld[€] \quad y = \leerfeld
\steil Welches Mietauto? Angebot 1: 29 € Miete pro Tag, 250 km insgesamt frei, jeder weitere km 0{,}40 €. Angebot 2: einmalig 110 € für eine Woche, jeder km 0{,}10 €. Jana mietet 3 Tage und fährt 380 km. Berechne die Gesamtkosten beider Angebote, entscheide, welches günstiger ist, und stelle für Angebot 2 eine Gleichung für die Kosten y (in €) bei x km für eine Woche auf. \hfill (P10 2023 OS) \\ Angebot 1: \leerfeld[€] \quad Angebot 2: \leerfeld[€] \quad y = \leerfeld
\teil Welches Busunternehmen? Nord verlangt 150 € Grundpreis und 1{,}20 € je km, Süd 1{,}80 € je km ohne Grundpreis. Das Ziel ist 120 km entfernt, der Bus fährt hin und zurück. Welches Unternehmen ist günstiger? Bleibt es dabei, wenn am Ziel noch ein Ausflug von 80 km dazukommt? Begründe. \hfill (P10 2016 OS)
\teil Welches Busunternehmen? Adler verlangt 90 € Grundpreis und 0{,}80 € je km, Biber 1{,}10 € je km ohne Grundpreis. Das Ziel ist 130 km entfernt, der Bus fährt hin und zurück. Welches Unternehmen ist günstiger? Bleibt es dabei, wenn am Ziel noch ein Ausflug von 60 km dazukommt? Begründe. \hfill (P10 2016 OS)
\teil Welche Gleichung gehört zu welchem Sachverhalt? Sachverhalt 1: E-Scooter: 1 € fürs Entsperren und 25 Cent je Minute. Sachverhalt 2: Schwimmbad: 25 € Jahresgebühr und 1 € je Besuch. y ist der Preis in €, x die Anzahl der Einheiten. Ordne jedem Sachverhalt eine der Gleichungen $y = 25x + 1$, $y = x + 25$, $y = 0{,}25x + 1$ zu. \hfill (P10 2021 OS) \\ Sachverhalt 1: \leerfeld \quad Sachverhalt 2: \leerfeld
\teil Welche Gleichung gehört zu welchem Sachverhalt? Sachverhalt 1: Parkhaus: 2 € bei der Einfahrt und 60 Cent je Stunde. Sachverhalt 2: Sportverein: 60 € Aufnahmegebühr und 2 € je Training. y ist der Preis in €, x die Anzahl der Einheiten. Ordne jedem Sachverhalt eine der Gleichungen $y = 60x + 2$, $y = 2x + 60$, $y = 0{,}60x + 2$ zu. \hfill (P10 2021 OS) \\ Sachverhalt 1: \leerfeld \quad Sachverhalt 2: \leerfeld
\steil Wann ist das Ziel erreicht? Lina hat 640 € auf dem Konto und zahlt jeden Monat 45 € ein. y ist das Guthaben in €, x die Anzahl der Monate. Stelle eine Gleichung für das Guthaben auf und berechne, nach wie vielen Monaten Lina mindestens 1\,500 € hat. \hfill (P10 2022 OS) \\ y = \leerfeld \quad nach \leerfeld Monaten
\steil Wann ist das Ziel erreicht? Jan hat 1\,150 € auf dem Konto und zahlt jeden Monat 70 € ein. y ist das Guthaben in €, x die Anzahl der Monate. Stelle eine Gleichung für das Guthaben auf und berechne, nach wie vielen Monaten Jan mindestens 2\,000 € hat. \hfill (P10 2022 OS) \\ y = \leerfeld \quad nach \leerfeld Monaten
\steil Funktionsgleichung: Ein Umzugsunternehmen berechnet 180 € Grundpreis und 1{,}20 € je km. Gib die Funktionsgleichung für den Preis y (in €) bei x km an. \hfill (P10 2016 OS) \\ y = \leerfeld
\steil Funktionsgleichung: Ein Reisebus kostet 250 € Grundpreis und 2{,}40 € je km. Gib die Funktionsgleichung für den Preis y (in €) bei x km an. \hfill (P10 2016 OS) \\ y = \leerfeld
\teil Kontostand nach einem Jahr? Ole hat 830 € auf dem Konto und zahlt jeden Monat 35 € ein. Wie viel Geld hat er nach einem Jahr? \hfill (P10 2022 OS) \\ \leerfeld[€]
\teil Bäume nach 6 Jahren? Eine Stadt hat 2\,340 Straßenbäume und pflanzt jedes Jahr 85 neue dazu. Wie viele Straßenbäume hat sie nach 6 Jahren, wenn keiner gefällt wird? \hfill (P10 2022 OS) \\ \leerfeld Bäume
\teil Werte ergänzen: Ein Wassertank wird gleichmäßig geleert. Die Tabelle zeigt die Wassermenge. Trage die Wassermenge zu Beginn und nach 40 Minuten ein. \hfill (P10 2021 OS)

\wertetabelle[,270,240,210,180,]{Zeit in min}{Wasser in l}{0,5,10,15,20,40}
\teil Werte ergänzen: Ein Akku wird gleichmäßig geladen. Die Tabelle zeigt den Ladestand. Trage den Ladestand zu Beginn und nach 90 Minuten ein. \hfill (P10 2021 OS)

\wertetabelle[,26,32,38,44,]{Zeit in min}{Ladung in \%}{0,10,20,30,40,90}
\teil Welcher Graph gehört zu welcher Firma? Firma Blitz verlangt 40 € Grundpreis und 1{,}20 € je km, Firma Flott 1{,}60 € je km ohne Grundpreis. Ordne die Graphen I und II den Firmen zu und begründe deine Wahl für Blitz. \hfill (P10 2016 OS) \\

\begin{ksys}[xmin=0,xmax=150,ymin=0,ymax=250,xstep=25,ystep=50,xlabel=x in km,ylabel=y in €,ablesen] \gerade{1.6}{0}{I} \gerade{1.2}{40}{II} \end{ksys}
Blitz: \leerfeld \quad Flott: \leerfeld
\teil Welcher Graph gehört zu welchem Verein? Verein Hoch verlangt 30 € Aufnahmegebühr und 5 € je Monat, Verein Weit 8 € je Monat ohne Aufnahmegebühr. Ordne die Graphen I und II den Vereinen zu und begründe deine Wahl für Hoch. \hfill (P10 2016 OS) \\

\begin{ksys}[xmin=0,xmax=12,ymin=0,ymax=100,xstep=1,ystep=10,xlabel=x in Monaten,ylabel=y in €,ablesen] \gerade{5}{30}{I} \gerade{8}{0}{II} \end{ksys}
Hoch: \leerfeld \quad Weit: \leerfeld
\teil Welches Diagramm passt? Angebot 1: 32 € Miete für einen Tag, 300 km frei, jeder weitere km 0{,}30 €. Angebot 2: einmalig 95 €, jeder km 0{,}12 €. Die Graphen A bis D zeigen Kosten y (in €) für einen Tag bei x km. Gib für jedes Angebot den passenden Graphen an. \hfill (P10 2023 OS) \\

\begin{ksys}[xmin=0,xmax=600,ymin=0,ymax=200,xstep=100,ystep=20,xlabel=x in km,ylabel=y in €,ablesen] \funktion{max(32,32+0.3*(\x-300))}{A} \gerade{0.12}{0}{B} \gerade{0.12}{95}{C} \gerade{0.3}{32}{D} \end{ksys}
Angebot 1: \leerfeld \quad Angebot 2: \leerfeld
\teil Welches Diagramm passt? Angebot 1: 25 € Miete für einen Tag, 200 km frei, jeder weitere km 0{,}40 €. Angebot 2: einmalig 80 €, jeder km 0{,}15 €. Die Graphen A bis D zeigen Kosten y (in €) für einen Tag bei x km. Gib für jedes Angebot den passenden Graphen an. \hfill (P10 2023 OS) \\

\begin{ksys}[xmin=0,xmax=500,ymin=0,ymax=160,xstep=100,ystep=20,xlabel=x in km,ylabel=y in €,ablesen] \gerade{0.4}{25}{A} \gerade{0.15}{80}{B} \gerade{0.15}{0}{C} \funktion{max(25,25+0.4*(\x-200))}{D} \end{ksys}
Angebot 1: \leerfeld \quad Angebot 2: \leerfeld
\end{teile}
\end{aufgabe}
